
\section{What I learned}

\paragraph{RELRO}

I learned about \href{https://www.redhat.com/en/blog/hardening-elf-binaries-using-relocation-read-only-relro}
{RELRO} through \texttt{checksec}; it consists in making the GOT table after
preliminary resolution of at the beginning of the execution.

\paragraph{Not following symlinks}

I learned about how not to follow symlinks using glibc functions; overall, I read a
whole lot of manual pages that I had never even gotten close to before.

\paragraph{File locks}

I had never before heard about file locks; now I know they exist, they are advisory
only and they come in \href{https://gavv.net/articles/file-locks/}{different flavours}
depending on the target platform; I ended up using BSD locks (\texttt{flock} from
\texttt{sys/file.h}).

\paragraph{Bundling a library statically}

For a long time I thought \textit{statically linking} was a binary characteristic of
a binary: either statically linked or dynamically linked. What I learned during this
project is that a binary can be linked statically against one library while
being dynamically linked against all the others.

This brings quite some advantages for e.g. a portable client where its library is
not commonly installed, without the need of bundling the whole glibc or openssl.

\paragraph{The security risk of RPATH/RUNPATH}

I learned about the security issue with \texttt{@rpath} (i.e. library loading from
a path relative to the execution one): running from a publicly writeable directory
will overwrite the dynamic loader path with one relative to said directory, thus
exposing the binary to loading a malicious library from there.

The usage of the build variable \texttt{\$ORIGIN} is preferrable, albeit still
not considered a best practice\cite{stackexchange-rpath-security}.

\section{Challenges encountered}

\paragraph{Error conventions}
A lot of concentration was keeping in adapting to the conventions of the various APIs
concerning error handling, one example being OpenSSL returning 1 for success and 0 for
failure, and another being \texttt{errno}-based libraries and those relying on
"custom-internal-error-storage" systems so to call them.

\paragraph{Authentication of the ciphertext}

My plan was initially to use AES-CBC as a cipher and compute an HMAC on the
ciphertext: while being cryptographically valid, this would have made the
implementation more complicated. For this reason I moved onto AES-GCM.

\paragraph{Cleanup paths within functions}

Duplicating the whole set of cleanup instructions for each failure path hinders
maintainability of the codebase through, precisely, code duplication.

At the same time, having a single cleanup section can lead to double free weaknesses
as well as \texttt{free} calls on uninitialized pointers.

This taught me that \texttt{goto} semantics aren't properly dead in C programming.

\section{LLM usage disclaimer}

I used GPT-6 for the following tasks...

\begin{itemize}
    \item Q\&A on the relevant topics
    \item refactoring of large functions, extraction of smaller ones
    \item implementation of the logging helper functions
    \item the entire implementation of the testing suite
    \item documentation of the testing suite in the report
    \item checking missing requirements
    \item bug hunting on the final draft of the program
    \item an abandoned attempt of porting the build system to support MacOS
    \item a very verbose yet comprehensive draft of the report on top
            of which I have written myself (credits to the IAIB programme 
            \href{https://gitlab.cs.taltech.ee/iaib/thesis-tex-estonian}{here}
            for the LaTeX template)
\end{itemize}

This list intends to be comprehensive of the work I have outsourced to generative AI
(and, as such, outsourced at all) to realize this project. Any omission, for how
unlikely I am to have missed something in the list above, is my genuine mistake.
